\unitchapter{2}{5}{System Software \& Operating Systems}{p2-05}

\unitmeta{Expected questions: 10--12 · Target: 10+ · Type: scheduling/paging/disk
numericals + Galvin concepts --- highly scoring}

\begin{sylbox}

\begin{itemize}
\tightlist
\item[$\square$]
  System software: assembly language, assemblers, compilers vs interpreters, loading, linking, relocation, macros, debuggers
\item[$\square$]
  OS basics: structure, services, system calls, boot
\item[$\square$]
  Process management: process states, PCB, IPC, synchronisation, critical section, Peterson, semaphores, classical problems
\item[$\square$]
  Threads: multithreading models, libraries, issues
\item[$\square$]
  CPU scheduling: criteria, FCFS, SJF, SRTF, priority, RR, multilevel queues, multiprocessor \& real-time scheduling
\item[$\square$]
  Deadlocks: conditions, RAG, prevention, avoidance (Banker\textquotesingle s), detection, recovery
\item[$\square$]
  Memory management: contiguous allocation, paging, segmentation, TLB, demand paging, page replacement, thrashing
\item[$\square$]
  Storage: disk structure, disk scheduling, RAID
\item[$\square$]
  File systems \& I/O: access methods, directories, allocation methods, free-space management, I/O subsystem
\item[$\square$]
  Security \& protection: access matrix, threats, authentication
\item[$\square$]
  Virtual machines; Linux \& Windows internals; distributed systems
\end{itemize}

\end{sylbox}

\needspace{95pt}

\hypertarget{p2-05--1-system-software}{%
\section{System Software}\label{p2-05--1-system-software}}

\diagram{figures/p2-05-00}

\begingroup\small

\par\medskip\noindent\hfil\begin{tabular}{@{}
  >{\raggedright\arraybackslash}p{(\columnwidth - 2\tabcolsep) * \real{0.1483}}
  >{\raggedright\arraybackslash}p{(\columnwidth - 2\tabcolsep) * \real{0.8517}}@{}}
\toprule\noalign{}
\begin{minipage}[b]{\linewidth}\raggedright
\textbf{Tool}
\end{minipage} & \begin{minipage}[b]{\linewidth}\raggedright
\textbf{Job}
\end{minipage} \\
\midrule\noalign{}
Assembler & Assembly → machine code. \textbf{Two-pass}: pass 1 builds symbol table (location counter), pass 2 generates code. One-pass needs \textbf{back-patching} for forward references \\
Compiler & Whole program → target code at once; faster execution \\
Interpreter & Statement by statement; slower, easier debugging \\
Linker & Resolves external references, combines object modules \\
Loader & Places program in memory; functions: \textbf{Allocation, Linking, Relocation, Loading} \\
Macro processor & Expands macros (text substitution, no call overhead); macro vs subroutine \\
Debugger & Breakpoints, single-step, watch variables \\
\bottomrule\noalign{}
\end{tabular}\hfil\null\par\medskip


\endgroup

\begin{itemize}
\tightlist
\item
  Loaders: \textbf{compile-and-go}, \textbf{absolute} (programmer does relocation), \textbf{relocating} (BSS loader), \textbf{direct-linking} loader, \textbf{dynamic loading} (load routine when called), \textbf{dynamic linking} (link at run time --- DLL, \texttt{.so}).
\item
  \textbf{Relocation}: adjusting address-sensitive instructions when program is loaded at a different address; relocation bits/table.
\item
  Assembler data structures: \textbf{Symbol table (ST)}, \textbf{Literal table (LT)}, \textbf{Opcode table (MOT/OPTAB)}, \textbf{Pseudo-op table (POT)}, Pool table.
\item
  Assembler directives (pseudo-ops): START, END, ORIGIN, EQU, LTORG, DC, DS.
\end{itemize}

\needspace{95pt}

\hypertarget{p2-05--2-os-basics}{%
\section{OS Basics}\label{p2-05--2-os-basics}}

\diagram{figures/p2-05-01}

\begingroup\small

\par\medskip\noindent\hfil\begin{tabular}{@{}
  >{\raggedright\arraybackslash}p{(\columnwidth - 2\tabcolsep) * \real{0.1894}}
  >{\raggedright\arraybackslash}p{(\columnwidth - 2\tabcolsep) * \real{0.7022}}@{}}
\toprule\noalign{}
\begin{minipage}[b]{\linewidth}\raggedright
\textbf{Structure}
\end{minipage} & \begin{minipage}[b]{\linewidth}\raggedright
\textbf{Example}
\end{minipage} \\
\midrule\noalign{}
Simple / monolithic & MS-DOS, early UNIX, Linux (monolithic + loadable modules) \\
Layered & THE OS (Dijkstra) \\
\textbf{Microkernel} & Mach, QNX, Minix --- minimal kernel, services in user space via message passing \\
Modular & Solaris, Linux LKMs \\
Hybrid & Windows NT, macOS (XNU) \\
Exokernel & Research --- exposes hardware \\
\bottomrule\noalign{}
\end{tabular}\hfil\null\par\medskip


\endgroup

\begin{itemize}
\tightlist
\item
  \textbf{Dual mode}: user mode / kernel mode (mode bit). System call → trap → kernel mode.
\item
  System call categories: process control (\texttt{fork}, \texttt{exec}, \texttt{exit}, \texttt{wait}), file management (\texttt{open}, \texttt{read}), device management (\texttt{ioctl}), information maintenance (\texttt{getpid}), communication (\texttt{pipe}, \texttt{shmget}), protection (\texttt{chmod}).
\item
  Parameter passing to system calls: registers, block/table in memory, stack.
\item
  Boot: power on → BIOS/UEFI firmware (POST) → bootstrap loader (MBR / GRUB) → kernel loaded → init/systemd.
\item
  OS types: batch, multiprogramming (maximise CPU utilisation), time-sharing (multitasking, response time), real-time (hard/soft), distributed, network, embedded.
\end{itemize}

\textbf{\texttt{fork()} puzzle}: \texttt{n} consecutive \texttt{fork()} calls create \textbf{2ⁿ − 1} child processes (2ⁿ total).

\begin{listingblock}{42pt}
\begin{Verbatim}[fontsize=\fontsize{9.0}{10.9}\selectfont]
fork(); fork(); fork();   // 8 processes total, 7 children
if (fork() && fork()) fork();   // count carefully: 4 processes total (parent + 3 children)
\end{Verbatim}
\end{listingblock}

\needspace{131pt}

\hypertarget{p2-05--3-process-management}{%
\section{Process Management}\label{p2-05--3-process-management}}

\needspace{95pt}

\hypertarget{p2-05--31-process-states}{%
\subsection{Process States}\label{p2-05--31-process-states}}

\diagram{figures/p2-05-02}

\begin{itemize}
\tightlist
\item
  \textbf{PCB} (Process Control Block): PID, state, PC, registers, scheduling info, memory info (page tables), accounting, I/O status.
\item
  Schedulers: \textbf{long-term} (job; controls degree of multiprogramming), \textbf{short-term} (CPU; very frequent), \textbf{medium-term} (swapping).
\item
  \textbf{Context switch} time is pure overhead. \textbf{Dispatcher latency}.
\item
  \textbf{Zombie}: terminated but parent hasn\textquotesingle t called \texttt{wait()}. \textbf{Orphan}: parent terminated (adopted by init).
\end{itemize}

\needspace{95pt}

\hypertarget{p2-05--32-ipc}{%
\subsection{IPC}\label{p2-05--32-ipc}}

\begin{itemize}
\tightlist
\item
  \textbf{Shared memory} (fast, needs synchronisation) vs \textbf{message passing} (send/receive; direct/indirect via mailboxes; blocking = synchronous, non-blocking = asynchronous).
\item
  Pipes (ordinary -- parent-child, unidirectional; named pipes/FIFOs), sockets (IP + port), RPC (stubs, marshalling), signals.
\end{itemize}

\needspace{229pt}

\hypertarget{p2-05--33-critical-section-problem}{%
\subsection{Critical Section Problem}\label{p2-05--33-critical-section-problem}}

Requirements: \textbf{Mutual exclusion}, \textbf{Progress}, \textbf{Bounded waiting}.

\textbf{Peterson\textquotesingle s solution (2 processes)}:

\begin{listingblock}{119pt}
\begin{Verbatim}[fontsize=\fontsize{9.0}{10.9}\selectfont]
// shared: bool flag[2]; int turn;
do {
    flag[i] = true;
    turn = j;                       // politely give turn to other
    while (flag[j] && turn == j);   // busy wait
    /* critical section */
    flag[i] = false;
    /* remainder section */
} while (true);
\end{Verbatim}
\end{listingblock}

Satisfies all three requirements (on sequentially consistent memory).

Hardware: \textbf{TestAndSet}, \textbf{CompareAndSwap} (atomic). Simple TAS lock gives ME + progress but not bounded waiting.

\needspace{92pt}

\hypertarget{p2-05--34-semaphores}{%
\subsection{Semaphores}\label{p2-05--34-semaphores}}

\begin{listingblock}{42pt}
\begin{Verbatim}[fontsize=\fontsize{9.0}{10.9}\selectfont]
wait(S)   { while (S <= 0); S--; }      // P(), down()
signal(S) { S++; }                      // V(), up()
\end{Verbatim}
\end{listingblock}

\begin{itemize}
\tightlist
\item
  \textbf{Binary semaphore} (mutex, 0/1) vs \textbf{counting semaphore}.
\item
  Blocking implementation uses a waiting queue --- negative value = number of waiting processes.
\item
  Counting semaphore initial value 10, then 6 P and 4 V → 10 − 6 + 4 = \textbf{8}.
\item
  Misuse → deadlock (wrong order of waits) or starvation. \textbf{Priority inversion} → solved by priority inheritance.
\item
  \textbf{Monitors}: high-level construct; condition variables with \texttt{wait()}/\texttt{signal()}.
\item
  \textbf{Spinlock} = busy-waiting lock (good for short waits on multiprocessors).
\end{itemize}

\needspace{130pt}

\hypertarget{p2-05--classical-problems}{%
\subsection{Classical problems}\label{p2-05--classical-problems}}

\begingroup\small

\par\medskip\noindent\hfil\begin{tabular}{@{}
  >{\raggedright\arraybackslash}p{(\columnwidth - 2\tabcolsep) * \real{0.3217}}
  >{\raggedright\arraybackslash}p{(\columnwidth - 2\tabcolsep) * \real{0.6783}}@{}}
\toprule\noalign{}
\begin{minipage}[b]{\linewidth}\raggedright
\textbf{Problem}
\end{minipage} & \begin{minipage}[b]{\linewidth}\raggedright
\textbf{Semaphores}
\end{minipage} \\
\midrule\noalign{}
Bounded buffer (producer--consumer) & \texttt{mutex\ =\ 1}, \texttt{empty\ =\ n}, \texttt{full\ =\ 0} \\
Readers--writers & \texttt{rw\_mutex\ =\ 1}, \texttt{mutex\ =\ 1}, \texttt{read\_count\ =\ 0} (first readers--writers: writers may starve) \\
Dining philosophers (5) & \texttt{chopstick{[}5{]}\ =\ 1}; deadlock if all pick left --- fixes: at most 4 at table, asymmetric pick, pick both atomically \\
Sleeping barber & customers, barbers, mutex \\
\bottomrule\noalign{}
\end{tabular}\hfil\null\par\medskip


\endgroup

\begin{listingblock}{97pt}
\begin{Verbatim}[fontsize=\fontsize{9.0}{10.9}\selectfont]
Producer:                       Consumer:
  wait(empty);                    wait(full);
  wait(mutex);                    wait(mutex);
  add item                        remove item
  signal(mutex);                  signal(mutex);
  signal(full);                   signal(empty);
(swapping the two waits in producer can cause DEADLOCK)
\end{Verbatim}
\end{listingblock}

\needspace{95pt}

\hypertarget{p2-05--4-threads}{%
\section{Threads}\label{p2-05--4-threads}}

\begin{itemize}
\tightlist
\item
  Thread = lightweight process: own \textbf{PC, registers, stack}; shares code, data, files with peers.
\item
  \textbf{User-level threads} (fast, kernel unaware; one blocks → all block in many-to-one) vs \textbf{kernel-level threads}.
\item
  Models: \textbf{Many-to-one} (Green threads), \textbf{One-to-one} (Linux, Windows), \textbf{Many-to-many}, two-level.
\item
  Libraries: POSIX Pthreads, Windows threads, Java threads. Implicit threading: thread pools, OpenMP, GCD, fork-join.
\item
  Issues: \texttt{fork()}/\texttt{exec()} semantics, signal handling, thread cancellation (asynchronous vs deferred), thread-local storage, scheduler activations (LWP, upcalls).
\item
  \textbf{Amdahl\textquotesingle s law} for multicore: speedup ≤ 1 / (S + (1−S)/N).
\end{itemize}

\needspace{114pt}

\hypertarget{p2-05--5-cpu-scheduling}{%
\section{CPU Scheduling}\label{p2-05--5-cpu-scheduling}}

\begin{listingblock}{64pt}
\begin{Verbatim}[fontsize=\fontsize{9.0}{10.9}\selectfont]
Turnaround time (TAT) = Completion − Arrival
Waiting time (WT)     = TAT − Burst
Response time         = First run − Arrival
Throughput            = #processes / total time
\end{Verbatim}
\end{listingblock}

\begingroup\small

\par\medskip\noindent\hfil\begin{tabular}{@{}
  >{\raggedright\arraybackslash}p{(\columnwidth - 4\tabcolsep) * \real{0.2374}}
  >{\raggedright\arraybackslash}p{(\columnwidth - 4\tabcolsep) * \real{0.2224}}
  >{\raggedright\arraybackslash}p{(\columnwidth - 4\tabcolsep) * \real{0.5402}}@{}}
\toprule\noalign{}
\begin{minipage}[b]{\linewidth}\raggedright
\textbf{Algorithm}
\end{minipage} & \begin{minipage}[b]{\linewidth}\raggedright
\textbf{Preemptive?}
\end{minipage} & \begin{minipage}[b]{\linewidth}\raggedright
\textbf{Notes}
\end{minipage} \\
\midrule\noalign{}
FCFS & No & Convoy effect \\
SJF & No & Optimal average waiting time (non-preemptive) \\
\textbf{SRTF} & Yes (preemptive SJF) & \textbf{Optimal} average WT overall; starvation \\
Priority & Both & Starvation → \textbf{aging} \\
\textbf{Round Robin} & Yes & Time quantum q; large q → FCFS; small q → high context switch overhead \\
Multilevel queue & --- & Fixed queues (foreground RR, background FCFS) \\
Multilevel feedback queue & Yes & Processes move between queues; most general \\
HRRN & No & Response ratio = (W + S)/S; favours short jobs, avoids starvation \\
\bottomrule\noalign{}
\end{tabular}\hfil\null\par\medskip


\endgroup

\needspace{130pt}

\hypertarget{p2-05--worked-example}{%
\subsection{Worked Example}\label{p2-05--worked-example}}

\begingroup\small

\par\medskip\noindent\hfil\begin{tabular}{@{}
  >{\raggedright\arraybackslash}p{(\columnwidth - 4\tabcolsep) * \real{0.0945}}
  >{\raggedright\arraybackslash}p{(\columnwidth - 4\tabcolsep) * \real{0.0887}}
  >{\raggedright\arraybackslash}p{(\columnwidth - 4\tabcolsep) * \real{0.0720}}@{}}
\toprule\noalign{}
\begin{minipage}[b]{\linewidth}\raggedright
\textbf{Process}
\end{minipage} & \begin{minipage}[b]{\linewidth}\raggedright
\textbf{Arrival}
\end{minipage} & \begin{minipage}[b]{\linewidth}\raggedright
\textbf{Burst}
\end{minipage} \\
\midrule\noalign{}
P1 & 0 & 8 \\
P2 & 1 & 4 \\
P3 & 2 & 9 \\
P4 & 3 & 5 \\
\bottomrule\noalign{}
\end{tabular}\hfil\null\par\medskip


\endgroup

\textbf{SRTF Gantt chart}:

\begin{listingblock}{42pt}
\begin{Verbatim}[fontsize=\fontsize{9.0}{10.9}\selectfont]
| P1 | P2      | P4        | P1             | P3                  |
0    1         5           10               17                    26
\end{Verbatim}
\end{listingblock}

\begingroup\small

\par\medskip\noindent\hfil\begin{tabular}{@{}
  >{\raggedright\arraybackslash}p{(\columnwidth - 6\tabcolsep) * \real{0.0458}}
  >{\raggedright\arraybackslash}p{(\columnwidth - 6\tabcolsep) * \real{0.0539}}
  >{\raggedright\arraybackslash}p{(\columnwidth - 6\tabcolsep) * \real{0.0719}}
  >{\raggedright\arraybackslash}p{(\columnwidth - 6\tabcolsep) * \real{0.0539}}@{}}
\toprule\noalign{}
\begin{minipage}[b]{\linewidth}\raggedright
\textbf{P}
\end{minipage} & \begin{minipage}[b]{\linewidth}\raggedright
\textbf{CT}
\end{minipage} & \begin{minipage}[b]{\linewidth}\raggedright
\textbf{TAT}
\end{minipage} & \begin{minipage}[b]{\linewidth}\raggedright
\textbf{WT}
\end{minipage} \\
\midrule\noalign{}
P1 & 17 & 17 & 9 \\
P2 & 5 & 4 & 0 \\
P3 & 26 & 24 & 15 \\
P4 & 10 & 7 & 2 \\
\bottomrule\noalign{}
\end{tabular}\hfil\null\par\medskip


\endgroup

Average WT = 26/4 = \textbf{6.5}.

\textbf{Round Robin (q = 4)} on same data:

\begin{listingblock}{42pt}
\begin{Verbatim}[fontsize=\fontsize{9.0}{10.9}\selectfont]
| P1 | P2 | P3 | P4 | P1 | P3 | P4 | P3 |
0    4    8    12   16   20   24   25   26
\end{Verbatim}
\end{listingblock}

CT: P1 = 20, P2 = 8, P3 = 26, P4 = 25 → WT: P1 = 12, P2 = 3, P3 = 15, P4 = 17 → average = \textbf{11.75}.

\needspace{95pt}

\hypertarget{p2-05--multiprocessor--real-time}{%
\subsection{Multiprocessor \& Real-time}\label{p2-05--multiprocessor--real-time}}

\begin{itemize}
\tightlist
\item
  Asymmetric vs \textbf{symmetric multiprocessing (SMP)}; processor affinity (soft/hard); load balancing (push/pull migration).
\item
  \textbf{Rate Monotonic} (static priority: shorter period → higher priority); schedulable if U ≤ n(2\^{}(1/n) − 1) (≈ 0.69 for large n).
\item
  \textbf{EDF} (Earliest Deadline First, dynamic): schedulable iff U ≤ 1.
\item
  Linux: \textbf{CFS} (Completely Fair Scheduler, red-black tree keyed by virtual runtime). Windows: 32-level priority-based preemptive.
\end{itemize}

\needspace{224pt}

\hypertarget{p2-05--6-deadlocks}{%
\section{Deadlocks}\label{p2-05--6-deadlocks}}

\needspace{188pt}

\hypertarget{p2-05--four-necessary-conditions-coffman}{%
\subsection{Four Necessary Conditions (Coffman)}\label{p2-05--four-necessary-conditions-coffman}}

\textbf{Mutual exclusion · Hold and wait · No preemption · Circular wait} --- all four must hold.

\begin{listingblock}{108pt}
\begin{Verbatim}[fontsize=\fontsize{9.0}{10.9}\selectfont]
Resource Allocation Graph
  P1 ──request──► [R1 •]──assigned──► P2
   ▲                                    │
   │                                    ▼
  [R2 •]◄─────────request───────────────┘
  (assigned to P1)
  Cycle with single-instance resources ⇒ DEADLOCK
  Cycle with multi-instance resources ⇒ deadlock POSSIBLE
\end{Verbatim}
\end{listingblock}

\begingroup\small

\par\medskip\noindent\hfil\begin{tabular}{@{}
  >{\raggedright\arraybackslash}p{(\columnwidth - 2\tabcolsep) * \real{0.1867}}
  >{\raggedright\arraybackslash}p{(\columnwidth - 2\tabcolsep) * \real{0.8133}}@{}}
\toprule\noalign{}
\begin{minipage}[b]{\linewidth}\raggedright
\textbf{Strategy}
\end{minipage} & \begin{minipage}[b]{\linewidth}\raggedright
\textbf{How}
\end{minipage} \\
\midrule\noalign{}
Prevention & Break one condition: e.g. request all at once (hold \& wait), allow preemption, \textbf{impose total order on resources} (circular wait) \\
Avoidance & \textbf{Banker\textquotesingle s algorithm} --- grant only if state remains safe \\
Detection \& recovery & Wait-for graph (single instance), detection algorithm; recover by killing processes / preemption \\
Ignorance & Ostrich algorithm (UNIX, Windows) \\
\bottomrule\noalign{}
\end{tabular}\hfil\null\par\medskip


\endgroup

\needspace{160pt}

\hypertarget{p2-05--bankers-algorithm--worked}{%
\subsection{Banker\textquotesingle s Algorithm --- Worked}\label{p2-05--bankers-algorithm--worked}}

5 processes, resources A(10), B(5), C(7).

\begingroup\small

\par\medskip\noindent\hfil\begin{tabular}{@{}
  >{\raggedright\arraybackslash}p{(\columnwidth - 6\tabcolsep) * \real{0.0458}}
  >{\raggedright\arraybackslash}p{(\columnwidth - 6\tabcolsep) * \real{0.2212}}
  >{\raggedright\arraybackslash}p{(\columnwidth - 6\tabcolsep) * \real{0.0667}}
  >{\raggedright\arraybackslash}p{(\columnwidth - 6\tabcolsep) * \real{0.2212}}@{}}
\toprule\noalign{}
\begin{minipage}[b]{\linewidth}\raggedright
\textbf{P}
\end{minipage} & \begin{minipage}[b]{\linewidth}\raggedright
\textbf{Allocation (A B C)}
\end{minipage} & \begin{minipage}[b]{\linewidth}\raggedright
\textbf{Max}
\end{minipage} & \begin{minipage}[b]{\linewidth}\raggedright
\textbf{Need = Max − Alloc}
\end{minipage} \\
\midrule\noalign{}
P0 & 0 1 0 & 7 5 3 & 7 4 3 \\
P1 & 2 0 0 & 3 2 2 & 1 2 2 \\
P2 & 3 0 2 & 9 0 2 & 6 0 0 \\
P3 & 2 1 1 & 2 2 2 & 0 1 1 \\
P4 & 0 0 2 & 4 3 3 & 4 3 1 \\
\bottomrule\noalign{}
\end{tabular}\hfil\null\par\medskip


\endgroup

Available = (10,5,7) − (7,2,5) = \textbf{(3,3,2)}.

\begin{listingblock}{86pt}
\begin{Verbatim}[fontsize=\fontsize{9.0}{10.9}\selectfont]
Work=(3,3,2): P1 need(1,2,2) ✓ → Work=(5,3,2)
              P3 need(0,1,1) ✓ → Work=(7,4,3)
              P4 need(4,3,1) ✓ → Work=(7,4,5)
              P0 need(7,4,3) ✓ → Work=(7,5,5)
              P2 need(6,0,0) ✓ → Work=(10,5,7)
Safe sequence: <P1, P3, P4, P0, P2>
\end{Verbatim}
\end{listingblock}

\textbf{Deadlock-free condition formula}: n processes each needing max k instances of a resource with R instances total → deadlock-free if \textbf{R ≥ n(k − 1) + 1}.
e.g. 3 processes each need 3 units → minimum R = 3×2 + 1 = \textbf{7}.

\needspace{131pt}

\hypertarget{p2-05--7-memory-management}{%
\section{Memory Management}\label{p2-05--7-memory-management}}

\needspace{95pt}

\hypertarget{p2-05--71-contiguous-allocation}{%
\subsection{Contiguous Allocation}\label{p2-05--71-contiguous-allocation}}

\begin{itemize}
\tightlist
\item
  Fixed partitions → \textbf{internal fragmentation}. Variable partitions → \textbf{external fragmentation} (fix with compaction).
\item
  Placement: \textbf{First fit}, \textbf{Best fit} (smallest adequate hole), \textbf{Worst fit} (largest), Next fit.
\item
  \textbf{50-percent rule}: with first fit, N allocated blocks → 0.5N blocks lost to fragmentation.
\end{itemize}

\needspace{95pt}

\hypertarget{p2-05--72-paging}{%
\subsection{Paging}\label{p2-05--72-paging}}

\diagram{figures/p2-05-03}

\begin{listingblock}{75pt}
\begin{Verbatim}[fontsize=\fontsize{9.0}{10.9}\selectfont]
Logical address bits = log₂(virtual space);  offset bits d = log₂(page size)
#pages = 2^(LA bits − d) ;  #frames = physical size / page size
Page table size = #pages × PTE size
EAT with TLB (hit ratio h, TLB t, memory m):
    EAT = h(t + m) + (1 − h)(t + 2m)
\end{Verbatim}
\end{listingblock}

\textbf{Worked}: 32-bit logical address, 4 KB pages, 4-byte PTE.

\begin{itemize}
\tightlist
\item
  Offset = 12 bits; pages = 2²⁰; page table = 2²⁰ × 4 B = \textbf{4 MB} per process (hence multilevel paging).
\item
  Two-level (10 \textbar{} 10 \textbar{} 12): outer table 1024 entries × 4 B = 4 KB = one page.
\item
  Inverted page table: one entry per \textbf{frame} (size ∝ physical memory), searched by hashing.
\end{itemize}

\needspace{95pt}

\hypertarget{p2-05--73-segmentation}{%
\subsection{Segmentation}\label{p2-05--73-segmentation}}

\begin{itemize}
\tightlist
\item
  Logical address = (segment number, offset); segment table has \textbf{base \& limit}; offset ≥ limit → trap.
\item
  No internal fragmentation, but external fragmentation. Segmentation with paging (Intel x86).
\end{itemize}

\needspace{92pt}

\hypertarget{p2-05--74-virtual-memory--demand-paging}{%
\subsection{Virtual Memory \& Demand Paging}\label{p2-05--74-virtual-memory--demand-paging}}

\begin{listingblock}{42pt}
\begin{Verbatim}[fontsize=\fontsize{9.0}{10.9}\selectfont]
EAT with page faults = (1 − p) × ma + p × page-fault service time
e.g. ma = 200 ns, service = 8 ms, p = 0.001 → EAT ≈ 200 + 0.001 × 8,000,000 = 8.2 µs
\end{Verbatim}
\end{listingblock}

Page fault steps: trap → check valid → find free frame → schedule disk read → update table → restart instruction.

\needspace{196pt}

\hypertarget{p2-05--75-page-replacement--worked}{%
\subsection{Page Replacement --- Worked}\label{p2-05--75-page-replacement--worked}}

Reference string: \textbf{7 0 1 2 0 3 0 4 2 3 0 3 2}, 3 frames.

\textbf{FIFO}

\begin{listingblock}{86pt}
\begin{Verbatim}[fontsize=\fontsize{9.0}{10.9}\selectfont]
Ref:  7  0  1  2  0  3  0  4  2  3  0  3  2
F1    7  7  7  2  2  2  2  4  4  4  0  0  0
F2       0  0  0  0  3  3  3  2  2  2  2  2
F3          1  1  1  1  0  0  0  3  3  3  3
Fault ✱  ✱  ✱  ✱     ✱  ✱  ✱  ✱  ✱  ✱
→ 10 faults
\end{Verbatim}
\end{listingblock}

\textbf{Optimal (replace page used farthest in future)} → \textbf{7 faults}.
\textbf{LRU (replace least recently used)} → \textbf{9 faults}.

\begin{itemize}
\tightlist
\item
  \textbf{Belady\textquotesingle s anomaly}: more frames → more faults; occurs in \textbf{FIFO}, not in stack algorithms (LRU, Optimal).
  Classic string: 1 2 3 4 1 2 5 1 2 3 4 5 → FIFO 3 frames: 9 faults, 4 frames: 10 faults.
\item
  Other algorithms: Second chance (clock), LFU, MFU, NRU (reference \& modify bits).
\item
  \textbf{Thrashing}: process spends more time paging than executing --- when Σ working sets \textgreater{} available frames; fix with \textbf{working-set model} or \textbf{page-fault frequency} control, reduce degree of multiprogramming.
\item
  Frame allocation: equal, proportional; global vs local replacement.
\item
  \textbf{Copy-on-write} for \texttt{fork()}. \textbf{Memory-mapped files}.
\end{itemize}

\needspace{196pt}

\hypertarget{p2-05--8-disk--storage}{%
\section{Disk \& Storage}\label{p2-05--8-disk--storage}}

\needspace{160pt}

\hypertarget{p2-05--81-disk-scheduling--worked}{%
\subsection{Disk Scheduling --- Worked}\label{p2-05--81-disk-scheduling--worked}}

Queue: \textbf{98, 183, 37, 122, 14, 124, 65, 67}; head at \textbf{53}; cylinders 0--199.

\begingroup\small

\par\medskip\noindent\hfil\begin{tabular}{@{}
  >{\raggedright\arraybackslash}p{(\columnwidth - 4\tabcolsep) * \real{0.2374}}
  >{\raggedright\arraybackslash}p{(\columnwidth - 4\tabcolsep) * \real{0.3794}}
  >{\raggedright\arraybackslash}p{(\columnwidth - 4\tabcolsep) * \real{0.3832}}@{}}
\toprule\noalign{}
\begin{minipage}[b]{\linewidth}\raggedright
\textbf{Algorithm}
\end{minipage} & \begin{minipage}[b]{\linewidth}\raggedright
\textbf{Order}
\end{minipage} & \begin{minipage}[b]{\linewidth}\raggedright
\textbf{Total head movement}
\end{minipage} \\
\midrule\noalign{}
FCFS & 53→\allowbreak{}98→\allowbreak{}183→\allowbreak{}37→\allowbreak{}122→\allowbreak{}14→\allowbreak{}124→\allowbreak{}65→\allowbreak{}67 & \textbf{640} \\
SSTF & 53→\allowbreak{}65→\allowbreak{}67→\allowbreak{}37→\allowbreak{}14→\allowbreak{}98→\allowbreak{}122→\allowbreak{}124→\allowbreak{}183 & \textbf{236} \\
SCAN (towards 0) & 53→\allowbreak{}37→\allowbreak{}14→\allowbreak{}0→\allowbreak{}65→\allowbreak{}67→\allowbreak{}98→\allowbreak{}122→\allowbreak{}124→\allowbreak{}183 & \textbf{236} \\
C-SCAN (towards 199) & 53→\allowbreak{}65→\allowbreak{}\ldots→\allowbreak{}183→\allowbreak{}199→\allowbreak{}0→\allowbreak{}14→\allowbreak{}37 & 146 + 199 + 37 = \textbf{382} (counting return jump) \\
LOOK (towards 0) & 53→\allowbreak{}37→\allowbreak{}14→\allowbreak{}65→\allowbreak{}\ldots→\allowbreak{}183 & 39 + 169 = \textbf{208} \\
C-LOOK (up) & 53→\allowbreak{}65→\allowbreak{}67→\allowbreak{}98→\allowbreak{}122→\allowbreak{}124→\allowbreak{}183→\allowbreak{}14→\allowbreak{}37 & 130 + 169 + 23 = \textbf{322} \\
\bottomrule\noalign{}
\end{tabular}\hfil\null\par\medskip


\endgroup

\begin{listingblock}{86pt}
\begin{Verbatim}[fontsize=\fontsize{9.0}{10.9}\selectfont]
SSTF head path (cylinder axis →)
0    14    37    53 65 67    98    122 124        183   199
                 ●──►──►                                    53→65→67
           ◄─────────────┘                                 67→37
      ◄────┘                                               37→14
      └──────────────────────►──────►──►──────────►         14→98→122→124→183
\end{Verbatim}
\end{listingblock}

\begin{itemize}
\tightlist
\item
  SSTF may cause \textbf{starvation}. SCAN = elevator. C-SCAN gives more uniform wait time.
\end{itemize}

\needspace{130pt}

\hypertarget{p2-05--82-raid}{%
\subsection{RAID}\label{p2-05--82-raid}}

\begingroup\small

\par\medskip\noindent\hfil\begin{tabular}{@{}
  >{\raggedright\arraybackslash}p{(\columnwidth - 6\tabcolsep) * \real{0.1006}}
  >{\raggedright\arraybackslash}p{(\columnwidth - 6\tabcolsep) * \real{0.3226}}
  >{\raggedright\arraybackslash}p{(\columnwidth - 6\tabcolsep) * \real{0.1139}}
  >{\raggedright\arraybackslash}p{(\columnwidth - 6\tabcolsep) * \real{0.2464}}@{}}
\toprule\noalign{}
\begin{minipage}[b]{\linewidth}\raggedright
\textbf{Level}
\end{minipage} & \begin{minipage}[b]{\linewidth}\raggedright
\textbf{Technique}
\end{minipage} & \begin{minipage}[b]{\linewidth}\raggedright
\textbf{Min disks}
\end{minipage} & \begin{minipage}[b]{\linewidth}\raggedright
\textbf{Fault tolerance}
\end{minipage} \\
\midrule\noalign{}
0 & Striping (block) & 2 & None \\
1 & Mirroring & 2 & 1 disk \\
2 & Bit striping + Hamming ECC & --- & 1 \\
3 & Byte striping + dedicated parity & 3 & 1 \\
4 & Block striping + dedicated parity & 3 & 1 (parity disk bottleneck) \\
\textbf{5} & Block striping + \textbf{distributed parity} & 3 & 1 \\
6 & Distributed double parity (P+Q) & 4 & 2 \\
10 (1+0) & Mirror then stripe & 4 & 1 per mirror pair \\
\bottomrule\noalign{}
\end{tabular}\hfil\null\par\medskip


\endgroup

\needspace{95pt}

\hypertarget{p2-05--9-file-systems--io}{%
\section{File Systems \& I/O}\label{p2-05--9-file-systems--io}}

\begin{itemize}
\tightlist
\item
  Access methods: sequential, direct (relative), indexed.
\item
  Directory structures: single-level → two-level → tree → \textbf{acyclic graph} (sharing via links) → general graph (cycles; needs garbage collection).
\item
  Hard link (same inode, same FS, no dirs) vs soft/symbolic link (path, can cross FS, can dangle).
\end{itemize}

\needspace{114pt}

\hypertarget{p2-05--allocation-methods}{%
\subsection{Allocation Methods}\label{p2-05--allocation-methods}}

\begin{listingblock}{64pt}
\begin{Verbatim}[fontsize=\fontsize{9.0}{10.9}\selectfont]
 Contiguous        Linked (FAT is a variant)      Indexed (inode)
 [A A A A]         A→A→A→A→nil                   index block → [b1,b2,b3,…]
 fast, direct      no ext. fragmentation          direct access, no ext. frag.
 ext. fragmentation  only sequential access        overhead of index block
\end{Verbatim}
\end{listingblock}

\textbf{UNIX inode max file size}: 12 direct + 1 single + 1 double + 1 triple indirect. Block 4 KB, pointer 4 B → 1024 pointers per block:

\begin{listingblock}{31pt}
\begin{Verbatim}[fontsize=\fontsize{9.0}{10.9}\selectfont]
(12 + 1024 + 1024² + 1024³) × 4 KB ≈ 4 TB
\end{Verbatim}
\end{listingblock}

\begin{itemize}
\tightlist
\item
  Free-space management: \textbf{bit vector} (1 bit per block), linked list, grouping, counting.
\item
  Directory implementation: linear list, hash table.
\item
  Mounting, file sharing, NFS. Journaling (log-structured) file systems for consistency.
\end{itemize}

\textbf{I/O}: polling, interrupts, DMA (see \protect\hyperlink{p2-02}{Unit 2}). Kernel I/O subsystem: scheduling, \textbf{buffering} (single, double, circular), \textbf{caching}, \textbf{spooling} (printer), error handling. Device drivers. Block vs character devices.

\needspace{95pt}

\hypertarget{p2-05--10-protection--security}{%
\section{Protection \& Security}\label{p2-05--10-protection--security}}

\begin{itemize}
\tightlist
\item
  \textbf{Access matrix}: rows = domains, columns = objects. Implementations: global table, \textbf{access control lists} (column-wise, per object), \textbf{capability lists} (row-wise, per domain), lock-key.
\item
  Revocation: immediate/delayed, selective/general, partial/total.
\item
  Principle of \textbf{least privilege}; need-to-know.
\item
  Program threats: Trojan horse, trap door (back door), logic bomb, stack/buffer overflow, virus. System/network threats: worms, port scanning, DoS/DDoS.
\item
  Authentication: passwords (salted hash), OTP, biometrics, multi-factor.
\item
  Cryptography --- see \protect\hyperlink{p2-09}{Unit 9}.
\end{itemize}

\needspace{95pt}

\hypertarget{p2-05--11-virtual-machines}{%
\section{Virtual Machines}\label{p2-05--11-virtual-machines}}

\begin{itemize}
\tightlist
\item
  \textbf{Type 1 (bare-metal) hypervisor}: VMware ESXi, Xen, Hyper-V, KVM. \textbf{Type 2 (hosted)}: VirtualBox, VMware Workstation.
\item
  Full virtualisation, para-virtualisation (guest OS modified --- Xen), hardware-assisted (Intel VT-x, AMD-V), containers (OS-level: Docker --- share host kernel), emulation (QEMU).
\item
  JVM, .NET CLR = application/process VMs.
\end{itemize}

\needspace{130pt}

\hypertarget{p2-05--12-linux--windows}{%
\section{Linux \& Windows}\label{p2-05--12-linux--windows}}

\begingroup\small

\par\medskip\noindent\hfil\begin{tabular}{@{}
  >{\raggedright\arraybackslash}p{(\columnwidth - 4\tabcolsep) * \real{0.1609}}
  >{\raggedright\arraybackslash}p{(\columnwidth - 4\tabcolsep) * \real{0.4315}}
  >{\raggedright\arraybackslash}p{(\columnwidth - 4\tabcolsep) * \real{0.4076}}@{}}
\toprule\noalign{}
\begin{minipage}[b]{\linewidth}\raggedright
\textbf{Feature}
\end{minipage} & \begin{minipage}[b]{\linewidth}\raggedright
\textbf{Linux}
\end{minipage} & \begin{minipage}[b]{\linewidth}\raggedright
\textbf{Windows}
\end{minipage} \\
\midrule\noalign{}
Kernel & Monolithic + loadable kernel modules & Hybrid (microkernel-inspired), HAL, executive \\
Process creation & \texttt{fork()} + \texttt{exec()} (clone for threads) & \texttt{CreateProcess()} \\
Scheduling & CFS (normal), real-time FIFO/RR & Priority-based preemptive, 32 levels \\
Memory & Buddy system + slab allocator, demand paging & Demand paging with clustering, working sets \\
File system & ext2/3/4 (inodes), VFS layer; \texttt{/proc} & NTFS (MFT, journaling, ACLs), FAT32, ReFS \\
IPC & pipes, signals, shared memory, message queues, sockets & ALPC, pipes, mailslots \\
Other & GPL; first released 1991 (Linus Torvalds) & Terminal services, fast user switching \\
\bottomrule\noalign{}
\end{tabular}\hfil\null\par\medskip


\endgroup

Linux process states: Running, Interruptible sleep, Uninterruptible sleep, Stopped, Zombie.
\textbf{Buddy system}: memory split into power-of-2 blocks; request 70 KB from 1 MB → gets 128 KB block.

\needspace{95pt}

\hypertarget{p2-05--13-distributed-systems}{%
\section{Distributed Systems}\label{p2-05--13-distributed-systems}}

\begin{itemize}
\tightlist
\item
  Network OS (users aware of machines; remote login, file transfer) vs \textbf{Distributed OS} (transparent: data, computation, process migration).
\item
  Topologies \& communication: naming, routing, connection strategies, contention.
\item
  Robustness: failure detection (heartbeat), reconfiguration, recovery.
\item
  Design issues: transparency, fault tolerance, scalability.
\item
  \textbf{Distributed File Systems}: NFS (stateless, v3), AFS (whole-file caching), GFS/HDFS. Naming transparency vs location independence; caching \& consistency (write-through, delayed-write, callback).
\item
  Clock synchronisation: \textbf{Lamport logical clocks} (happened-before →), vector clocks, Cristian\textquotesingle s, Berkeley algorithms.
\item
  Mutual exclusion: centralised, Ricart--Agrawala (2(n−1) messages), token ring. Election: \textbf{Bully}, Ring.
\end{itemize}

\needspace{166pt}

\hypertarget{p2-05--14-deeper-dive--more-scheduling-semaphore-traces-multilevel-paging--file-numericals}{%
\section{Deeper Dive --- More Scheduling, Semaphore Traces, Multilevel Paging \& File Numericals}\label{p2-05--14-deeper-dive--more-scheduling-semaphore-traces-multilevel-paging--file-numericals}}

\needspace{130pt}

\hypertarget{p2-05--141-priority--hrrn-scheduling--worked}{%
\subsection{Priority \& HRRN Scheduling --- Worked}\label{p2-05--141-priority--hrrn-scheduling--worked}}

\begingroup\small

\par\medskip\noindent\hfil\begin{tabular}{@{}
  >{\raggedright\arraybackslash}p{(\columnwidth - 6\tabcolsep) * \real{0.0979}}
  >{\raggedright\arraybackslash}p{(\columnwidth - 6\tabcolsep) * \real{0.0919}}
  >{\raggedright\arraybackslash}p{(\columnwidth - 6\tabcolsep) * \real{0.0745}}
  >{\raggedright\arraybackslash}p{(\columnwidth - 6\tabcolsep) * \real{0.2719}}@{}}
\toprule\noalign{}
\begin{minipage}[b]{\linewidth}\raggedright
\textbf{Process}
\end{minipage} & \begin{minipage}[b]{\linewidth}\raggedright
\textbf{Arrival}
\end{minipage} & \begin{minipage}[b]{\linewidth}\raggedright
\textbf{Burst}
\end{minipage} & \begin{minipage}[b]{\linewidth}\raggedright
\textbf{Priority (lower = higher)}
\end{minipage} \\
\midrule\noalign{}
P1 & 0 & 10 & 3 \\
P2 & 0 & 1 & 1 \\
P3 & 0 & 2 & 4 \\
P4 & 0 & 1 & 5 \\
P5 & 0 & 5 & 2 \\
\bottomrule\noalign{}
\end{tabular}\hfil\null\par\medskip


\endgroup

\textbf{Non-preemptive priority}:

\begin{listingblock}{42pt}
\begin{Verbatim}[fontsize=\fontsize{9.0}{10.9}\selectfont]
| P2 | P5    | P1         | P3 | P4 |
0    1       6            16   18   19
\end{Verbatim}
\end{listingblock}

Waiting times: P1 6, P2 0, P3 16, P4 18, P5 1 → average = 41/5 = \textbf{8.2}.

\textbf{HRRN} (non-preemptive; response ratio RR = (W + S)/S), processes A(0, 3), B(2, 6), C(4, 4), D(6, 5), E(8, 2):

\begin{listingblock}{86pt}
\begin{Verbatim}[fontsize=\fontsize{9.0}{10.9}\selectfont]
t = 0: only A → run A to 3
t = 3: only B arrived → run B to 9
t = 9: C: (5 + 4)/4 = 2.25   D: (3 + 5)/5 = 1.6   E: (1 + 2)/2 = 1.5   → run C to 13
t = 13: D: (7 + 5)/5 = 2.4   E: (5 + 2)/2 = 3.5                        → run E to 15
t = 15: run D to 20
Order: A B C E D
\end{Verbatim}
\end{listingblock}

\needspace{103pt}

\hypertarget{p2-05--142-semaphore-tracing}{%
\subsection{Semaphore Tracing}\label{p2-05--142-semaphore-tracing}}

\begin{listingblock}{53pt}
\begin{Verbatim}[fontsize=\fontsize{9.0}{10.9}\selectfont]
semaphore S = 1, T = 0;
P1: wait(S); print("A"); signal(T);
P2: wait(T); print("B"); signal(S);
\end{Verbatim}
\end{listingblock}

Whatever the scheduling, P2 blocks on T until P1 signals → output \textbf{"AB"} (and the pair can repeat as A B A B \ldots{} if looped). This is the standard \textbf{ordering} use of semaphores (initialise to 0 to force "happens-after").

\textbf{Deadlock with wrong order}:

\begin{listingblock}{42pt}
\begin{Verbatim}[fontsize=\fontsize{9.0}{10.9}\selectfont]
P0: wait(S); wait(Q); …        P1: wait(Q); wait(S); …     (S = Q = 1)
P0 gets S, P1 gets Q, each waits for the other → deadlock
\end{Verbatim}
\end{listingblock}

\needspace{144pt}

\hypertarget{p2-05--143-multilevel-paging-numericals}{%
\subsection{Multilevel Paging Numericals}\label{p2-05--143-multilevel-paging-numericals}}

48-bit virtual address, 4 KB pages, 8-byte PTEs.

\begin{listingblock}{64pt}
\begin{Verbatim}[fontsize=\fontsize{9.0}{10.9}\selectfont]
Offset = 12 bits → page-number bits = 36
Entries per page-table page = 4096 / 8 = 512 = 2⁹ → each level indexes 9 bits
Levels needed = ⌈36 / 9⌉ = 4  (this is x86-64 4-level paging: 9 | 9 | 9 | 9 | 12)
Memory accesses per reference without TLB = 4 (tables) + 1 (data) = 5
\end{Verbatim}
\end{listingblock}

\textbf{Effective access time with TLB and 2-level paging}: TLB 10 ns, memory 100 ns, hit ratio 90\%:
EAT = 0.9 × (10 + 100) + 0.1 × (10 + 3 × 100) = 99 + 31 = \textbf{130 ns}.

\needspace{130pt}

\hypertarget{p2-05--144-segmentation--address-translation}{%
\subsection{Segmentation --- Address Translation}\label{p2-05--144-segmentation--address-translation}}

\begingroup\small

\par\medskip\noindent\hfil\begin{tabular}{@{}
  >{\raggedright\arraybackslash}p{(\columnwidth - 4\tabcolsep) * \real{0.1022}}
  >{\raggedright\arraybackslash}p{(\columnwidth - 4\tabcolsep) * \real{0.0655}}
  >{\raggedright\arraybackslash}p{(\columnwidth - 4\tabcolsep) * \real{0.0777}}@{}}
\toprule\noalign{}
\begin{minipage}[b]{\linewidth}\raggedright
\textbf{Segment}
\end{minipage} & \begin{minipage}[b]{\linewidth}\raggedright
\textbf{Base}
\end{minipage} & \begin{minipage}[b]{\linewidth}\raggedright
\textbf{Limit}
\end{minipage} \\
\midrule\noalign{}
0 & 219 & 600 \\
1 & 2300 & 14 \\
2 & 90 & 100 \\
3 & 1327 & 580 \\
\bottomrule\noalign{}
\end{tabular}\hfil\null\par\medskip


\endgroup

\begin{itemize}
\tightlist
\item
  (0, 430) → 430 \textless{} 600 → physical \textbf{649}
\item
  (1, 10) → \textbf{2310}
\item
  (2, 500) → 500 ≥ 100 → \textbf{segmentation fault (trap)}
\item
  (3, 400) → \textbf{1727}
\end{itemize}

\needspace{125pt}

\hypertarget{p2-05--145-page-replacement-with-4-frames-belady-check}{%
\subsection{Page Replacement with 4 Frames (Belady check)}\label{p2-05--145-page-replacement-with-4-frames-belady-check}}

Reference: 1 2 3 4 1 2 5 1 2 3 4 5

\begin{itemize}
\tightlist
\item
  FIFO, 3 frames → \textbf{9} faults; 4 frames → \textbf{10} faults (anomaly).
\item
  LRU, 4 frames → \textbf{8} faults; Optimal, 4 frames → \textbf{6} faults.
\end{itemize}

\needspace{111pt}

\hypertarget{p2-05--146-file-system-numericals}{%
\subsection{File-System Numericals}\label{p2-05--146-file-system-numericals}}

\textbf{Max file size with UNIX inode}: block 1 KB, pointer 4 B → 256 pointers/block; 10 direct, 1 single, 1 double, 1 triple:

\begin{listingblock}{31pt}
\begin{Verbatim}[fontsize=\fontsize{9.0}{10.9}\selectfont]
(10 + 256 + 256² + 256³) × 1 KB = (10 + 256 + 65,536 + 16,777,216) KB ≈ 16 GB
\end{Verbatim}
\end{listingblock}

\textbf{FAT size}: 4 GB disk, 4 KB clusters, 32-bit entries → 2²⁰ clusters × 4 B = \textbf{4 MB} FAT.

\textbf{Disk access}: 7200 RPM, average seek 8 ms, transfer rate 100 MB/s, read 4 KB block:
seek 8 + latency 4.17 + transfer 0.04 ≈ \textbf{12.2 ms}.

\needspace{144pt}

\hypertarget{p2-05--147-buddy-system--worked}{%
\subsection{Buddy System --- Worked}\label{p2-05--147-buddy-system--worked}}

1 MB block; requests A = 70 KB, B = 35 KB, C = 80 KB.

\begin{listingblock}{64pt}
\begin{Verbatim}[fontsize=\fontsize{9.0}{10.9}\selectfont]
1024 → 512 + 512 → 256 + 256 → 128 + 128          A gets 128 KB (wastes 58 KB internal)
B (35 KB): split a 128 → 64 + 64                   B gets 64 KB
C (80 KB): needs 128 → split the second 256 → 128 + 128, C gets 128 KB
Release: buddies of equal size and adjacent addresses coalesce back
\end{Verbatim}
\end{listingblock}

\needspace{147pt}

\hypertarget{p2-05--148-process-synchronisation-hardware}{%
\subsection{Process Synchronisation Hardware}\label{p2-05--148-process-synchronisation-hardware}}

\begin{listingblock}{97pt}
\begin{Verbatim}[fontsize=\fontsize{9.0}{10.9}\selectfont]
// Test-and-Set lock
bool lock = false;
do {
    while (test_and_set(&lock));   // atomically: old = lock; lock = true; return old
    /* critical section */
    lock = false;
} while (true);
\end{Verbatim}
\end{listingblock}

Gives mutual exclusion and progress, but \textbf{not bounded waiting} (a process can be overtaken indefinitely); the bounded-waiting version uses a \texttt{waiting{[}{]}} array.

\needspace{125pt}

\hypertarget{p2-05--previous-year-questions-pyq-pattern}{%
\section{Previous Year Questions (PYQ pattern)}\label{p2-05--previous-year-questions-pyq-pattern}}

\textbf{System software}

\begin{enumerate}
\def\labelenumi{\arabic{enumi}.}
\tightlist
\item
  A two-pass assembler\textquotesingle s first pass: \textbf{builds the symbol table (assigns addresses)}
\item
  Forward reference problem in one-pass assembler is solved by: \textbf{back-patching}
\item
  Which loader function adjusts address-dependent locations? \textbf{Relocation}
\item
  Dynamic linking is performed at: \textbf{run time}
\item
  Macro expansion is done by: \textbf{macro processor / preprocessor} (before assembly/compilation)
\end{enumerate}

\textbf{Processes \& threads}

\begin{enumerate}
\def\labelenumi{\arabic{enumi}.}
\setcounter{enumi}{5}
\tightlist
\item
  Number of child processes created by \texttt{for(i=0;i\textless{}n;i++)\ fork();}: \textbf{2ⁿ − 1}
\item
  Which is NOT shared by threads of a process? \textbf{Stack (and registers)}
\item
  Degree of multiprogramming is controlled by: \textbf{long-term scheduler}
\item
  A process whose parent has not called wait(): \textbf{zombie}
\item
  Which multithreading model blocks the whole process on a blocking system call? \textbf{Many-to-one}
\end{enumerate}

\textbf{Synchronisation}

\begin{enumerate}
\def\labelenumi{\arabic{enumi}.}
\setcounter{enumi}{10}
\tightlist
\item
  Critical section solution requirements: \textbf{mutual exclusion, progress, bounded waiting}
\item
  Semaphore S = 7; 20 P and 15 V operations performed: final = 7 − 20 + 15 = \textbf{2}
\item
  Producer-consumer with bounded buffer of size n: initial values mutex = 1, empty = \textbf{n}, full = \textbf{0}
\item
  Priority inversion is solved by: \textbf{priority inheritance protocol}
\item
  Peterson\textquotesingle s solution is for: \textbf{two processes}
\end{enumerate}

\textbf{Scheduling}

\begin{enumerate}
\def\labelenumi{\arabic{enumi}.}
\setcounter{enumi}{15}
\tightlist
\item
  Algorithm with minimum average waiting time: \textbf{SJF / SRTF}
\item
  Round robin with very large quantum becomes: \textbf{FCFS}
\item
  Starvation in priority scheduling is solved by: \textbf{aging}
\item
  Convoy effect occurs in: \textbf{FCFS}
\item
  Processes P1(0,10), P2(1,1), P3(2,2) with SRTF: average waiting time = P1: 3, P2: 0, P3: 0 → \textbf{1}
\item
  Rate-monotonic schedulability bound for 2 tasks: 2(√2 − 1) ≈ \textbf{0.828}
\end{enumerate}

\textbf{Deadlock}

\begin{enumerate}
\def\labelenumi{\arabic{enumi}.}
\setcounter{enumi}{21}
\tightlist
\item
  Necessary conditions for deadlock: \textbf{mutual exclusion, hold \& wait, no preemption, circular wait}
\item
  Banker\textquotesingle s algorithm is for deadlock: \textbf{avoidance}
\item
  Ordering resources numerically prevents: \textbf{circular wait}
\item
  4 processes each need 2 tape drives; minimum drives to ensure no deadlock: 4×1 + 1 = \textbf{5}
\item
  Cycle in RAG with single-instance resources implies: \textbf{deadlock}
\end{enumerate}

\textbf{Memory}

\begin{enumerate}
\def\labelenumi{\arabic{enumi}.}
\setcounter{enumi}{26}
\tightlist
\item
  Belady\textquotesingle s anomaly occurs in: \textbf{FIFO}
\item
  Page size 4 KB, logical address 32 bits → number of pages: \textbf{2²⁰}
\item
  TLB 20 ns, memory 100 ns, hit 80\%: EAT = 0.8×120 + 0.2×220 = \textbf{140 ns}
\item
  Thrashing is caused by: \textbf{high degree of multiprogramming / insufficient frames}
\item
  Internal fragmentation occurs in: \textbf{paging (fixed-size)}; external in: \textbf{segmentation / variable partitions}
\item
  Best fit, memory holes 100K, 500K, 200K, 300K, 600K; request 212K → \textbf{300K hole}
\item
  LRU faults for 1,2,3,4,1,2,5,1,2,3,4,5 with 3 frames: \textbf{10}
\item
  Optimal page replacement is not implementable because: \textbf{future references are unknown}
\item
  Working set model is used to prevent: \textbf{thrashing}
\end{enumerate}

\textbf{Disk / file}

\begin{enumerate}
\def\labelenumi{\arabic{enumi}.}
\setcounter{enumi}{35}
\tightlist
\item
  Disk scheduling that may cause starvation: \textbf{SSTF}
\item
  Elevator algorithm: \textbf{SCAN}
\item
  RAID level with distributed parity: \textbf{RAID 5}
\item
  RAID with mirroring: \textbf{RAID 1}
\item
  Allocation method that supports direct access and has no external fragmentation: \textbf{Indexed}
\item
  FAT uses which allocation method? \textbf{Linked (table variant)}
\item
  Bit vector for 1 TB disk with 4 KB blocks: 2²⁸ bits = \textbf{32 MB}
\end{enumerate}

\textbf{Security / VM / distributed}

\begin{enumerate}
\def\labelenumi{\arabic{enumi}.}
\setcounter{enumi}{42}
\tightlist
\item
  Access control lists correspond to: \textbf{columns of the access matrix}
\item
  Capability list corresponds to: \textbf{rows (domains)}
\item
  A program that appears useful but performs malicious actions: \textbf{Trojan horse}
\item
  VMware ESXi is a: \textbf{Type 1 hypervisor}
\item
  Lamport clocks capture: \textbf{happened-before ordering (logical time)}
\item
  Linux kernel is: \textbf{monolithic with loadable modules}
\item
  Windows file system with MFT: \textbf{NTFS}
\end{enumerate}

\textbf{More practice questions}

\begin{enumerate}
\def\labelenumi{\arabic{enumi}.}
\setcounter{enumi}{49}
\tightlist
\item
  HRRN favours: \textbf{short jobs, while preventing starvation of long jobs (waiting time raises the ratio)}
\item
  Response ratio of a process with waiting time 9 and burst 3: \textbf{4}
\item
  Semaphore initialised to 0 is used for: \textbf{ordering / signalling between processes}
\item
  32-bit virtual address, 4 KB pages, 4-byte PTEs, single-level: page table size = \textbf{4 MB}
\item
  Levels of paging for 48-bit VA, 4 KB pages, 8 B PTEs (page-sized tables): \textbf{4}
\item
  Segment (2, 500) with limit 100: \textbf{trap / segmentation fault}
\item
  LRU with 4 frames on 1 2 3 4 1 2 5 1 2 3 4 5: \textbf{8} faults
\item
  Optimal with 4 frames on the same string: \textbf{6} faults
\item
  Buddy system allocation for a 70 KB request: \textbf{128 KB} block
\item
  Test-and-set spin lock fails to guarantee: \textbf{bounded waiting}
\item
  Average rotational latency at 15,000 RPM: \textbf{2 ms}
\end{enumerate}

\begin{revbox}

\begin{itemize}
\tightlist
\item
  n forks → 2ⁿ processes · Threads share code/data/files, not stack/registers
\item
  SRTF optimal WT · RR big q = FCFS · aging fixes starvation
\item
  Deadlock-free: R ≥ n(k−1)+1 · Banker = avoidance
\item
  EAT(TLB) = h(t+m) + (1−h)(t+2m) · Belady only FIFO
\item
  SSTF starvation · SCAN elevator · RAID 5 distributed parity
\end{itemize}

\end{revbox}
